% discussion.tex -- Discussion

Across eight Pfam families spanning a ten-fold range of sizes ($K{=}37$--$420$) and a seven-fold range of sequence lengths ($L{=}23$--$262$), SA generated novel sequences that preserved family-level amino acid statistics and were predicted by independent structure predictors to adopt the target-family fold. It did so using only the seed alignment, with no training, external data, or GPU resources. None of the learned baselines tested here achieved this combination. Profile HMMs and the MSA Transformer generated sequences with nearest-neighbor identity below the empirical within-family nearest-neighbor envelope. In five families (RRM, SH3, Kunitz, PDZ, and Pkinase), this identity fell below the minimum observed among natural sequences (SI Appendix, Table~\ref{tab:nn-band}). Their high novelty was therefore accompanied by sequence-identity drift rather than family-consistent exploration. EvoDiff required millions of pretraining sequences and hours of compute, and degraded on the longest family. Among the tested baselines, no method simultaneously stayed within the family identity range, maintained low composition divergence, and produced genuine novelty. That a training-free method matched systems built on orders of magnitude more data indicated that the statistical structure of protein families can support generation without iterative parameter optimization, backpropagation, or external pretraining, provided the generative framework respects the geometry of the family's sequence space. Beyond matching their output, SA avoids the pretraining corpora, GPU resources, and large parameter counts that place deep generative methods out of reach for the small, single-laboratory families that make up most of Pfam, requiring only a seed alignment and a standard workstation. The consensus-with-noise and column-permuted alignment controls supported the conclusion that this was not only a consequence of marginal statistics: SA produced a different output profile from isotropic perturbation of the consensus, and exploited inter-position correlations that column-independent models could not capture.

Sequences generated by stochastic attention were predicted to adopt the canonical family fold. In six of eight families, their TM-scores to the reference structure modestly exceeded those of the stored natural sequences under ESMFold (AlphaFold2 reproduced the pattern across all eight families, the two predictors' TM-scores correlating at $r{=}0.999$; SI Appendix, Fig.~\ref{fig:esmfold-af2-concordance}). The pLDDT scores were indistinguishable from stored sequences in six of eight families, supporting the conclusion that the sampler did not sacrifice predicted folding confidence for diversity, and the single exception (SH3) was resolved by AlphaFold2 cross-validation. Two explanations for the elevated TM-scores merit consideration. First, predictor bias: structure predictors trained on evolutionary data may assign higher confidence to consensus-proximal sequences. Second, genuine structural convergence: the Boltzmann ensemble averages over lineage-specific deviations, potentially producing sequences that encode the shared fold more cleanly than individual natural sequences. A direct test of the first explanation is whether predicted structural quality tracks consensus proximity within each family. The correlation was weak and inconsistent (pooled Spearman $\rho{=}0.17$ for TM-score and $0.23$ for pLDDT, each accounting for under $4\%$ of the variance, and negative in three families; SI Appendix, Section~\ref{si:consensus-reg}). This argues against predictor bias as the sole driver, though it does not exclude an absolute preference for family-typical sequences shared by both the sampler and the predictor. We therefore interpret the elevated TM-scores cautiously and do not read them as evidence that SA-generated sequences fold better than natural ones. The two explanations are not mutually exclusive. The decisive computational control, folding the consensus-with-noise ensemble through the same structure predictors, was not performed here and remains for future work. Definitive attribution requires experimental structure determination. The DMS validation provided partial evidence for the latter. If elevated TM-scores reflected only predictor bias, one would not necessarily expect SA substitutions to be enriched for experimentally tolerated mutations. Yet the position-matched enrichment of $1.09$--$1.21\times$ in two independent DMS datasets suggested that the generated substitutions are biased toward locally function-compatible changes. However, all computational validators used here (structure predictors, ESM2-650M, and DMS-derived fitness scores) were trained on evolutionary data, so high scores for family-like sequences are expected. Experimental characterization would provide definitive confirmation.

The WW domain scaling study showed that generation quality was stable down to $K{=}20$ sequences, with KL divergence below $0.02$ and novelty above $0.47$, and replication across three additional families supported the same trend. This robustness arose because the energy landscape is defined entirely by the stored patterns, with no parameters that could overfit to a small training set. The entropy inflection criterion adapted the operating temperature to the memory matrix at each scale, and $\beta^*$ could be predicted from PCA dimensionality alone ($R^2{=}0.97$), enabling fully automatic operation from a seed alignment. The $\sqrt{d}$ scaling of $\beta^*$ is consistent with concentration of measure on the unit sphere: random-probe similarities have variance $1/d$, so the inverse temperature required to resolve $\mathcal{O}(1/\sqrt{d})$ similarity differences grows as $\sqrt{d}$. The reduced coefficient ($0.28$ vs.\ ${\sim}1.6$ for random patterns) reflected self-similarity anchoring in the stored-pattern probes. Each stored pattern has unit inner product with itself but lower similarity to other patterns. The softmax therefore concentrated on the self-pattern at lower $\beta$ than a mean-field model predicts (SI Appendix, Section~\ref{si:beta-star-theory}). By contrast, learned generative models face a parameter-to-data ratio problem when $K < 100$: VAEs and diffusion models have thousands to millions of parameters that cannot be constrained by a few dozen sequences. This regime encompasses most Pfam families.

The connection between attention and Boltzmann sampling extends beyond proteins: any set of examples representable as vectors defines an energy landscape via the modern Hopfield energy, and Langevin dynamics on this landscape provides a generative model with no training. Two conditions determine whether the approach produces useful output. First, the stored examples must be structured enough that the energy landscape has well-defined basins rather than a featureless plateau. Protein families satisfy this because evolutionary constraint creates correlated patterns of conservation and variation. Second, the examples must be diverse enough that the sampler does not trivially retrieve: if all stored patterns are nearly identical, the distribution collapses and generation reduces to consensus with noise. The eight families spanned $22$--$37\%$ mean pairwise identity, diverse enough to avoid consensus collapse yet, through their correlated conservation and variation, structured enough to form well-defined basins. Classical Hopfield networks were limited to storing $\mathcal{O}(d)$ patterns before catastrophic interference~\cite{hopfieldNeuralNetworksPhysical1982,amitStoringInfiniteNumbers1985}. The modern continuous Hopfield energy lifts this limit, storing a number of well-separated patterns that grows exponentially with $d$~\cite{ramsauerHopfieldNetworksAll2021}. This is an advantage for scarce-data generation: with $K \ll d$ stored patterns the family lies below capacity, so the energy landscape is well-separated and the sampler can explore between memories without interference. The same approach may apply to other biological sequence families where data is scarce but examples are structured. Examples include antibody repertoires, where clonotypes have only a handful of somatic variants; RNA families with few seed sequences; and small-molecule libraries with tens of analogs per series. The Hopfield energy also admits a conditioning mechanism: adjusting pattern multiplicities via a scalar bias on the attention logits shifts generation toward a user-specified functional subset, enabling property-directed generation (e.g., toward binding activity) without retraining~\cite{varnerConditioningProtein2026}. The analytic score function makes the sampler amenable to the convergence theory of Langevin dynamics~\cite{durmusNonasymptoticAnalysisUnadjusted2017,dalalyanTheoreticalGuaranteesApproximate2017}. In the strongly regularized, log-concave regime ($\beta\sigma_{\max}^2 < 2$, with $\sigma_{\max}$ the largest singular value of the memory matrix), these results yield formal guarantees on sample quality that are absent from most neural generative models. The generation temperatures used here lie outside that regime ($\beta\sigma_{\max}^2 \gg 2$; SI Appendix, Section~\ref{si:convergence}), where the energy is non-convex and convergence is supported by the empirical mixing diagnostics and initialization-insensitivity analysis rather than by these bounds.

Several open questions remain for SA, each pointing to a direction for future development. The most significant concerns the fixed alignment representation: SA requires a pre-computed MSA and a fixed alignment length $L$, and cannot generate insertions, deletions, or variable-length sequences. This constraint is inherited from the one-hot-plus-PCA encoding, which maps each sequence to a vector of fixed dimension $20L$. Any change in $L$ would alter the representation space itself. In practice, this means the method is best suited to compact, well-aligned domains rather than multi-domain proteins or families with frequent large insertions. Extending stochastic attention to variable-length generation would likely require replacing PCA with a length-agnostic embedding while preserving the closed-form score function, an open direction for future work. The PCA projection itself discards information that may be relevant for capturing higher-order correlations, though a sensitivity analysis confirmed robustness across the 75\%--99\% retained variance range (SI Appendix, Table~\ref{tab:pca-sensitivity}). Relatedly, SA preserved pairwise residue covariation in seven of eight families but largely lost it in Pkinase ($r{=}0.05$), the longest and most data-starved family, where generation was consensus-proximal (SI Appendix, Table~\ref{tab:mi}). Covariation fidelity therefore degrades when the data-to-length ratio is low, and our covariation claims are restricted to the compact families. Further, the eight families studied here are all well-structured, globular domains with abundant structural and functional annotation. This choice was necessary for validation but introduced a selection bias toward the class of proteins that structure predictors handle best. Performance on intrinsically disordered regions, transmembrane domains, multi-domain proteins, or families with low sequence identity remains untested, and the method's reliance on PCA may be limiting for disordered families where sequence variation is less structured. Finally, our multi-chain protocol does not guarantee coverage of all modes in the Boltzmann distribution, particularly for large or diverse families, and wet-lab characterization of selected generated sequences would establish the method's utility for protein engineering applications.
